% ID: FIS-CIN-MUA-001
% Titolo: Spazio e tempo di frenata di un'automobile
% Disciplina: Fisica
% Area: Cinematica
% Argomento: Moto rettilineo uniformemente accelerato
% Sottoargomento: Frenata con accelerazione costante
% Classe: Seconda liceo scientifico
% Difficolta: 3
% Tipo: problema_numerico
% Risultato: a = -4,0 m/s^2; t = 5,0 s
% Tempo_stimato: 8 min
% Competenze: conversione delle unita, uso delle equazioni del moto, interpretazione del segno
% Prerequisiti: velocita, accelerazione, moto uniformemente accelerato
% Tag: fisica, cinematica, moto_accelerato, frenata, automobile
% Fonte: esempio_originale
% Autore: OpenAI Codex
% Licenza: CC-BY-SA-4.0

\begin{esercizio}
Un'automobile procede in rettilineo alla velocita di
\(72\,\mathrm{km/h}\). Il conducente frena e l'auto si arresta dopo
\(50\,\mathrm{m}\), con accelerazione costante.

Determina l'accelerazione dell'automobile durante la frenata e il tempo
necessario per fermarsi.
\end{esercizio}

\begin{soluzione}
Convertiamo prima la velocita iniziale:
\[
72\,\mathrm{km/h}=
\frac{72}{3{,}6}\,\mathrm{m/s}=20\,\mathrm{m/s}.
\]

Usiamo la relazione:
\[
v^2=v_0^2+2as.
\]

Poiche alla fine l'automobile e ferma, \(v=0\). Quindi:
\[
0=(20)^2+2a\cdot 50.
\]

Da cui:
\[
a=-\frac{400}{100}=-4{,}0\,\mathrm{m/s^2}.
\]

Il segno meno indica che l'accelerazione e opposta al verso della velocita.
Per trovare il tempo usiamo:
\[
v=v_0+at.
\]

\[
0=20-4{,}0t.
\]

\[
t=\frac{20}{4{,}0}=5{,}0\,\mathrm{s}.
\]
\end{soluzione}
